\documentclass{article}

\usepackage{amsmath,amssymb}
\usepackage{xurl}
\usepackage[hidelinks]{hyperref}
\setlength{\emergencystretch}{2em}

\input{command}

\title{Top-index scope of the spectral n-positivity criterion (Wolf Prop.~3.2)}
\author{}
\date{2026-06-24}

\begin{document}
\maketitle
\gapnote{scope-restriction}{resolved}

\section{Source statement}

Let $\tau\in\mathcal M_{D'}\otimes\mathcal M_D$ be Hermitian, with normalized
eigenvectors $\phi_i$ and spectral data
\begin{align}
    \tau&=\sum_i\nu_i\ketbra{\phi_i}{\phi_i},
    \label{eq:wolf_prop_3_2_top_index_scope_spectral_decomposition}\\
    \rho_i&=\tr_2\!\left(\ketbra{\phi_i}{\phi_i}\right).
    \label{eq:wolf_prop_3_2_top_index_scope_reduced_density}
\end{align}
Here $\rho_i$ acts on the first factor $\C^{D'}$. Let $\nu_0$ be the smallest
positive eigenvalue of $\tau$, let $\nu$ be its largest eigenvalue, and let
$\lVert\rho_i\rVert_{(n)}$ denote its Ky--Fan $n$-norm.

Proposition~3.2 of Ref.~\cite{Wolf2012Quantum} gives
\begin{align}
    \inf_\psi\bra{\psi}\tau\ket{\psi}
        &\geq \nu_0+
        \sum_{i:\nu_i\leq0}
        (\nu_i-\nu_0)\lVert\rho_i\rVert_{(n)}.
    \label{eq:wolf_prop_3_2_top_index_scope_lower_bound}
\end{align}
This is Equation~(3.7) of Ref.~\cite{Wolf2012Quantum}. If a single eigenvalue
$\nu_-$ is non-positive and $\rho_-$ is the reduced density matrix of its
eigenvector, the proposition also gives
\begin{align}
    \inf_\psi\bra{\psi}\tau\ket{\psi}
        &\leq \nu+(\nu_--\nu)\lVert\rho_-\rVert_{(n)}.
    \label{eq:wolf_prop_3_2_top_index_scope_upper_bound}
\end{align}
This is Equation~(3.8) of Ref.~\cite{Wolf2012Quantum}. The local source is
\path{Notes/WolfNoteTexSource/ch03_positive_not_completely.tex}, lines
153--179.

Because $\rho_i$ acts on $\C^{D'}$, its Ky--Fan norm is naturally indexed by
$1\leq n\leq D'$. Proposition~3.1 of Ref.~\cite{Wolf2012Quantum} restricts
the positivity index to $1\leq n\leq D$, as recorded in the same local source
at lines 90--92. Consequently, the endpoint $n=D'$ belongs to the source range
exactly when $D'\leq D$. If $D'>D$, that endpoint extends beyond the source
range. The formal statements nevertheless hold for every $n\geq1$, and hence
cover every index used by the source.

\section{Formalized statement below the top index}

For $1\leq n<D'$, the formal per-vector theorems prove that every normalized
vector $\psi$ of Schmidt rank at most $n$ satisfies the lower bound, and that
some such vector satisfies the upper bound.\footnote{Formal declarations:
    \leanid{Matrix.IsHermitian.spectral_lower_bound} and
    \leanid{Matrix.IsHermitian.exists_le_spectral_upper_bound} in
    \doclink{QICLean/Channel/NPositivitySpectralCriterion}.}

The infimum formulations use the set of normalized vectors of Schmidt rank at
most $n$.\footnote{Formal declarations:
    \leanid{Matrix.IsHermitian.le_sInf_schmidtRankLEExpectations} for
    Equation~(3.7) of Ref.~\cite{Wolf2012Quantum},
    \leanid{Matrix.IsHermitian.sInf_schmidtRankLEExpectations_le} for
    Equation~(3.8) of Ref.~\cite{Wolf2012Quantum}, and
    \leanid{Matrix.IsHermitian.sInf_schmidtRankLEExpectations_top} for the top
    index, all in
    \doclink{QICLean/Channel/NPositivitySpectralCriterion}.}
The interpretation of the source phrase ``vectors of Schmidt rank $n$'' is
discussed in
\path{docs/paper-gaps/wolf_prop_3_2_schmidt_rank_reading.tex}.

\section{The top index}

At $n=D'$, every normalized vector satisfies the Schmidt-rank constraint, and
each reduced density has full Ky--Fan sum
\begin{align}
    \lVert\rho_i\rVert_{(D')}
        &=\tr(\rho_i)
    \label{eq:wolf_prop_3_2_top_index_scope_ky_fan_trace}\\
        &=1.
    \label{eq:wolf_prop_3_2_top_index_scope_ky_fan_one}
\end{align}
Let $\nu_{\min}$ denote the smallest eigenvalue of $\tau$. The spectral
expansion gives, for every normalized $\psi$,
\begin{align}
    \bra{\psi}\tau\ket{\psi}
        &=\sum_i\nu_i
            \left|\braket{\phi_i}{\psi}\right|^2,
    \label{eq:wolf_prop_3_2_top_index_scope_rayleigh_expansion}\\
    \sum_i\left|\braket{\phi_i}{\psi}\right|^2
        &=1,
    \label{eq:wolf_prop_3_2_top_index_scope_parseval}\\
    \bra{\psi}\tau\ket{\psi}
        &\geq\nu_{\min}.
    \label{eq:wolf_prop_3_2_top_index_scope_rayleigh_lower_bound}
\end{align}
Evaluating at a normalized least-eigenvalue eigenvector gives
\begin{align}
    \inf_{\lVert\psi\rVert=1}\bra{\psi}\tau\ket{\psi}
        &=\nu_{\min}.
    \label{eq:wolf_prop_3_2_top_index_scope_minimum_eigenvalue}
\end{align}

At the top index, the right-hand side of the lower bound in Equation~(3.7) of
Ref.~\cite{Wolf2012Quantum} becomes
\begin{align}
    \nu_0+
    \sum_{i:\nu_i\leq0}
        (\nu_i-\nu_0)\lVert\rho_i\rVert_{(D')}
        &=
    \nu_0+\sum_{i:\nu_i\leq0}(\nu_i-\nu_0).
    \label{eq:wolf_prop_3_2_top_index_scope_lower_endpoint_value}
\end{align}
This quantity is at most $\nu_{\min}$ and can be strictly smaller. For the
spectrum $(-2,-1,3)$ with $\nu_0=3$,
\begin{align}
    \nu_0+\sum_{i:\nu_i\leq0}(\nu_i-\nu_0)
        &=3+(-2-3)+(-1-3)
    \label{eq:wolf_prop_3_2_top_index_scope_example_calculation}\\
        &=-6,
    \label{eq:wolf_prop_3_2_top_index_scope_example_value}\\
    \nu_{\min}&=-2.
    \label{eq:wolf_prop_3_2_top_index_scope_example_minimum}
\end{align}
Thus the top-index lower bound follows from
(\ref{eq:wolf_prop_3_2_top_index_scope_rayleigh_lower_bound}), but it is not
identical to the Rayleigh inequality.

Under the hypothesis of Equation~(3.8) of Ref.~\cite{Wolf2012Quantum},
$\nu_-$ is the unique non-positive eigenvalue and therefore equals
$\nu_{\min}$. Using
(\ref{eq:wolf_prop_3_2_top_index_scope_ky_fan_one}), its right-hand side is
\begin{align}
    \nu+(\nu_--\nu)\lVert\rho_-\rVert_{(D')}
        &=\nu+(\nu_--\nu)
    \label{eq:wolf_prop_3_2_top_index_scope_upper_endpoint_first}\\
        &=\nu_-
    \label{eq:wolf_prop_3_2_top_index_scope_upper_endpoint_value}\\
        &=\nu_{\min}.
    \label{eq:wolf_prop_3_2_top_index_scope_upper_endpoint_minimum}
\end{align}
A normalized least-eigenvalue eigenvector realizes this value. The endpoint
proof therefore uses the Rayleigh expansion and Parseval's identity, not the
Ky--Fan machinery required for $n<D'$.

The top-index minimum is formalized by
\leanid{Matrix.IsHermitian.spectral_lower_bound_top} and
\leanid{Matrix.IsHermitian.exists_spectral_lower_bound_top}. The top-index
infimum form of the lower bound is
\leanid{Matrix.IsHermitian.le_sInf_schmidtRankLEExpectations_top}; the upper
bound is the corresponding half of
\leanid{Matrix.IsHermitian.sInf_schmidtRankLEExpectations_top}.

\section{Classification}

The endpoint is resolved. Equations~(3.7) and~(3.8) of
Ref.~\cite{Wolf2012Quantum} are formalized at every index $n\geq1$, and hence
throughout the source range $1\leq n\leq D$, regardless of the relation between
$D$ and $D'$. The declarations
\leanid{Matrix.IsHermitian.spectral_lower_bound} and
\leanid{Matrix.IsHermitian.exists_le_spectral_upper_bound} cover
$1\leq n<D'$. The declarations
\leanid{Matrix.IsHermitian.spectral_lower_bound_top} and
\leanid{Matrix.IsHermitian.exists_spectral_lower_bound_top} cover $n\geq D'$ by
the endpoint argument in
(\ref{eq:wolf_prop_3_2_top_index_scope_rayleigh_expansion})--%
(\ref{eq:wolf_prop_3_2_top_index_scope_minimum_eigenvalue}). The infimum form
has the same coverage, with
\leanid{Matrix.IsHermitian.le_sInf_schmidtRankLEExpectations_top} supplying the
top-index lower bound. No hypothesis absent from
Ref.~\cite{Wolf2012Quantum} is used.

\bibliographystyle{alpha}
\bibliography{references}

\end{document}
